\begin{abstract}
This paper investigates minimal, resource-bounded runtime verification for safety-relevant automotive ECU telemetry. A server-side replay pipeline was implemented using LARVA monitors and benchmarked under two regimes: baseline (clean trace) and abnormal (injected violations), with runtime verification toggled on and off on the same endpoint. Four properties were evaluated: voltage out-of-range under running RPM, coolant over-temperature, throttle spike, and load--MAP mismatch. On the unmodified generated monitor (one warmup and five measured runs per mode), baseline overhead was 6.43\% (125.8 ms RV on vs 118.2 ms RV off), while abnormal overhead was 3927.05\% (4824.4 ms vs 119.8 ms), showing that monitor cost is strongly state-dependent. Injected violations produced high RV-on flag counts while RV-off flags remained zero. Threshold tuning reduced baseline false positives for load--MAP mismatch without removing abnormal detections. A controlled three-variant experiment isolated the dominant overhead source: relative to that stock monitor, abnormal overhead fell to 683.79\% with logging disabled and to 9.39\% after removing bad-state stacktrace construction, with CPU-time measurements following the same pattern. These results show that diagnostic generation, rather than guard evaluation, dominates violation-heavy runtime cost.
\end{abstract}